\chapter{От обычной симметрии к алгебре Грассмана}

\lectureoverview{Суперсимметрия является симметрией, поэтому мы начнем с напоминания о том, что означает это слово в квантовой механике. Физическое преобразование представляется в пространстве Гильберта унитарным оператором; непрерывные преобразования порождаются бесконечно малыми, а их коммутаторы кодируют способ композиции этих преобразований. Подлинная симметрия сохраняет энергию, что эквивалентно коммутированию с гамильтонианом. Далее суперсимметрия изменяет одно предположение, общее для всех знакомых примеров: ее генераторы переводят бозоны в фермионы, а фермионы в бозоны. Построение таких генераторов естественным образом приводит к операторам рождения и уничтожения, антикоммутаторам и конечной алгебре Грассмана, в которой нечетные величины в квадрате равны нулю.}

\section{Симметрия связывает разные конфигурации}

Симметрия — не просто внешняя правильность. Это операция, которая связывает одну физическую конфигурацию с другой. Лево-правая зеркальная симметрия, например, связывает две руки. Если бы человек и окружающие законы обладали точной лево-правой симметрией, соответствующие свойства двух рук совпадали бы. В более общем случае объекты или состояния, связанные точной симметрией, должны иметь одинаковые физические величины, которые эта симметрия сохраняет. В физике частиц это часто означает равные массы; в атомной физике — равные энергии состояний, принадлежащих одному и тому же вращательному мультиплету.

Наша цель — суперсимметрия, но сначала полезно вспомнить непрерывные симметрии, такие как вращения и трансляции. Вопрос о том, является ли суперсимметрия «непрерывной» в обычном смысле, уже вызывает затруднение. Ее бесконечно малые параметры не являются обычными коммутирующими вещественными числами. Ее лучше рассматривать как градуированное расширение непрерывной симметрии, и введенная ниже арифметика Грассмана делает это утверждение достаточно точным для наших нынешних целей. \editorialnote{В современной терминологии суперсимметрия описывается как непрерывная градуированная симметрия. Это уточнение объясняет предупреждение в лекции о том, что обычное противопоставление непрерывного и дискретного плохо применимо к суперсимметрии.}

\section{Унитарные операции в пространстве Гильберта}

Квантовая система представляется вектором состояния \(\ket{\psi}\). Физическая операция симметрии представляется оператором \(U\), действующим на это состояние:
\begin{equation}
\ket{\psi}\longmapsto U\ket{\psi}.
\end{equation}
Для пространственного вращения ось и угол определяют \(U\). Если состояние содержит частицы в нескольких положениях, преобразованный кет представляет их все в повернутых положениях. Трансляция действует аналогично: она заменяет состояние другим, в котором смещена вся конфигурация.

\begin{classroomqa}
\audiencequestion{Вектор состояния принадлежит пространству Гильберта, часто бесконечномерному. Мы вращаем пространство Гильберта или обычное трехмерное пространство? \lecturetimestamp{00:04:20}}
\lecturerresponse{Оба описания относятся к одной и той же операции, но на разных уровнях. Унитарное преобразование действует на вектор в пространстве Гильберта и представляет физическое вращение конфигурации в обычном пространстве. Это не буквальное трехмерное вращение абстрактного пространства Гильберта; это представление в нем пространственного вращения.}
\end{classroomqa}

Существенное условие — унитарность:
\begin{equation}
U^\dagger U=UU^\dagger=\mathbf 1,
\qquad
U^{-1}=U^\dagger.
\end{equation}
Это сохраняет скалярные произведения и, следовательно, вероятности:
\begin{equation}
\braket{U\psi}{U\psi}
=\bra{\psi}U^\dagger U\ket{\psi}
=\braket{\psi}{\psi}.
\end{equation}

\begin{figure}[H]
\centering
\includegraphics[width=0.56\textwidth]{lecture_04_figure_02.png}
\lectureframe{00:05:56.420}
\caption{Ранний переход на доске от преобразованного состояния к условию унитарности. Кет и произведение с эрмитовым сопряжением видны, но лектор закрывает их правые части; полные уравнения приведены в выводе ниже.}
\end{figure}

Строго говоря, теорема Вигнера допускает также антиунитарные симметрии, такие как обращение времени. Преобразования, непрерывно связанные с тождественным преобразованием, которые и рассматриваются здесь, представляются унитарными операторами. \editorialnote{Возможность антиунитарности — стандартное уточнение; она не меняет обсуждения в лекции вращений, трансляций или генераторов суперсимметрии.}

\section{Бесконечно малые преобразования и генераторы}

Ничего не делать представляется как \(U=\mathbf 1\), а не \(U=0\). Достаточно малое преобразование отличается от тождественного членом, линейным по его малому параметру. Для вращения на малый угол \(\epsilon\),
\begin{equation}
U(\epsilon)=\mathbf 1+i\epsilon L+O(\epsilon^2).
\end{equation}
Унитарность в первом порядке требует
\begin{align}
U^\dagger U
&=\left(\mathbf 1-i\epsilon L^\dagger\right)
  \left(\mathbf 1+i\epsilon L\right)+O(\epsilon^2)\\
&=\mathbf 1+i\epsilon(L-L^\dagger)+O(\epsilon^2),
\end{align}
и, следовательно,
\begin{equation}
L^\dagger=L.
\end{equation}
Оператор \(L\) является генератором бесконечно малого преобразования.

Для вращений нужны три генератора, по одному для каждой независимой оси:
\begin{equation}
L_x,\qquad L_y,\qquad L_z.
\end{equation}
Вращение вокруг произвольной оси \(\widehat{\bm n}\) порождается оператором
\begin{equation}
L_{\widehat{\bm n}}=\widehat{\bm n}\cdot\bm L.
\end{equation}

\begin{classroomqa}
\audiencequestion{Можно ли сложить генераторы для осей \(x\), \(y\) и \(z\), чтобы получить вращение вокруг другой оси? \lecturetimestamp{00:07:50}}
\lecturerresponse{Да. Их линейные комбинации порождают вращения вокруг промежуточных осей. Трех независимых компонент достаточно, чтобы построить любое бесконечно малое пространственное вращение.}
\end{classroomqa}

Для вращений этими генераторами являются компоненты углового момента. Такое обозначение полезно и в более общем случае, поскольку любая непрерывная симметрия локально описывается через свои генераторы.

\section{Коммутатор как таблица умножения}

Генераторы вращений удовлетворяют
\begin{align}
[L_x,L_y]&=iL_z,&
[L_y,L_z]&=iL_x,&
[L_z,L_x]&=iL_y,
\end{align}
или, в компактной записи,
\begin{equation}
[L_i,L_j]=i\epsilon_{ijk}L_k.
\end{equation}
Символ Леви-Чивиты обращается в нуль, если два индекса совпадают, и равен \(+1\) или \(-1\) для четных или нечетных перестановок \(1,2,3\). Важный структурный факт состоит в замкнутости: коммутатор двух генераторов является еще одним генератором.

Почему именно коммутатор должен описывать геометрию вращений? Выполните малое вращение вокруг \(x\), затем вокруг \(y\), отмените вращение вокруг \(x\) и, наконец, отмените вращение вокруг \(y\). Операции не сокращаются, потому что вращения вокруг разных осей не коммутируют.

\begin{classroomqa}
\audiencequestion{Можно ли наглядно представить эту некоммутативность на примере крена, тангажа и рыскания самолета? \lecturetimestamp{00:12:36}}
\lecturerresponse{Да. Последовательные вращения вокруг осей, связанных с телом, сразу геометрически показывают, что изменение порядка меняет конечную ориентацию. Алгебра ниже является бесконечно малой версией этого хорошо известного факта.}
\end{classroomqa}

Пусть
\begin{equation}
U_x(\epsilon)=\mathbf 1+i\epsilon L_x,
\qquad
U_y(\delta)=\mathbf 1+i\delta L_y.
\end{equation}
Поскольку операторы действуют на кеты справа налево, указанная последовательность представляется как
\begin{equation}
\mathcal C
=U_y^{-1}(\delta)U_x^{-1}(\epsilon)
 U_y(\delta)U_x(\epsilon).
\end{equation}

\begin{classroomqa}
\audiencequestion{Почему операции отмены можно записывать через обратные операторы, а не через эрмитово-сопряженные? \lecturetimestamp{00:13:25}}
\lecturerresponse{Потому что преобразования унитарны: \(U^\dagger=U^{-1}\). Для бесконечно малого преобразования взятие обратного просто меняет знак у члена первого порядка, \(U_x^{-1}=\mathbf 1-i\epsilon L_x+O(\epsilon^2)\).}
\end{classroomqa}

Перемножение четырех множителей на доске превратилось в оживленное упражнение: все члены, линейные по \(\epsilon\) или \(\delta\), сокращаются, аудитория следила за смешанными членами, а временная путаница с коэффициентом 2 показала, насколько легко потерять порядок операторов. Аккуратный результат второго порядка таков:
\begin{align}
\mathcal C
&=\mathbf 1+\epsilon\delta[L_x,L_y]
  +O(\epsilon^2\delta,\epsilon\delta^2)\\
&=\mathbf 1+i\epsilon\delta L_z
  +O(\epsilon^2\delta,\epsilon\delta^2).
\end{align}
Итак, замкнутая последовательность оставляет малое вращение вокруг \(z\) на угол, пропорциональный \(\epsilon\delta\). \editorialnote{Выписанное произведение фиксирует порядок и коэффициент после исследовательского расчета в аудитории. При изменении соглашения о последовательности знак меняется на противоположный, но вывод о том, что остаточное преобразование определяется коммутатором, остается тем же.}

\begin{classroomqa}
\audiencequestion{Сделает ли изменение порядка обратных операций произведение тождественным? \lecturetimestamp{00:25:23}}
\lecturerresponse{Если каждое преобразование поставить непосредственно рядом с его собственным обратным, пары сократятся. В нетривиальной коммутаторной последовательности эти пары намеренно разделены; то, что она не схлопывается, и измеряет зависимость от порядка.}
\end{classroomqa}

Более длинные произведения действительно становятся сложнее, но никакого нового принципа в них нет.

\begin{classroomqa}
\audiencequestion{Что происходит, когда генераторов больше трех или когда имеется много малых преобразований? \lecturetimestamp{00:25:34}}
\lecturerresponse{Коммутаторов становится больше, как и технического учета, но знание замкнутой алгебры позволяет определить, как комбинируется сколь угодно большое число бесконечно малых преобразований. Конечные преобразования можно собрать из многих малых.}
\end{classroomqa}

В этом смысле алгебра коммутаторов есть таблица умножения непрерывной группы, записанная вблизи тождества.

\section{Трансляции завершают пространственный пример}

Пространственные трансляции порождаются импульсом. В принятом в лекции соглашении о знаках,
\begin{equation}
\begin{aligned}
T_x(a)&=\mathbf 1+iaP_x+O(a^2),\\
T_y(b)&=\mathbf 1+ibP_y+O(b^2),
\end{aligned}
\end{equation}
где \(a\) и \(b\) — расстояния, а не углы.

Две трансляции коммутируют. Сместите точку по \(x\), затем по \(y\), отмените трансляцию по \(x\) и отмените трансляцию по \(y\); точка вернется в исходное положение:
\begin{equation}
[P_i,P_j]=0.
\end{equation}

\begin{classroomqa}
\audiencequestion{Требуются ли силы, чтобы начать и остановить движение при переносе или вращении физического объекта? \lecturetimestamp{00:28:45}}
\lecturerresponse{Определяется не эта операция. Преобразование симметрии сопоставляет исходное состояние с математическим состоянием, в котором вся конфигурация целиком смещена или повернута. Оно не описывает динамическую траекторию, ускорение или работу, совершаемую при перемещении установки.}
\end{classroomqa}

Поворот и перенос не коммутируют. Сместите точку, поверните её вокруг начала координат, перенесите обратно вдоль исходной оси и отмените поворот; точка не вернётся туда, откуда начала. Смешанная алгебра имеет вид
\begin{equation}
[L_i,P_j]=i\epsilon_{ijk}P_k.
\end{equation}
Это соотношение означает, что сам импульс преобразуется как вектор при поворотах. Точное упражнение с индексами здесь менее важно, чем общий вывод: преобразования евклидова пространства закодированы в коммутаторах между их генераторами поворотов и переносов.

\section{Симметрия должна сохранять энергию}

Теперь перейдём от композиции преобразований к динамике. Если \(U\) действительно является симметрией системы, его применение не может изменить энергию. Пусть \(\ket{E}\) — собственное состояние энергии:
\begin{equation}
H\ket{E}=E\ket{E}.
\end{equation}
Символ \(E\) справа — это число; \(\ket{E}\) — вектор, помеченный этим собственным значением.

\begin{figure}[H]
\centering
\begin{minipage}[t]{0.43\textwidth}
\centering
\includegraphics[width=\linewidth]{lecture_04_figure_03.png}
\lectureframe{00:34:38.239}
\end{minipage}
\hfill
\begin{minipage}[t]{0.50\textwidth}
\centering
\includegraphics[width=\linewidth]{lecture_04_figure_04.png}
\lectureframe{00:36:22.376}
\end{minipage}
\caption{Преобразованный кет \(U\ket{E}\) и определяющее уравнение на собственное значение \(H\ket{E}=E\ket{E}\), выделенные на последовательных кадрах на доске.}
\end{figure}

Если \(U\) — симметрия, то \(U\ket{E}\) имеет то же собственное значение:
\begin{equation}
HU\ket{E}=EU\ket{E}.
\end{equation}
Так как \(E\) — скаляр, он коммутирует с \(U\), и можно использовать исходное уравнение на собственные значения:
\begin{align}
HU\ket{E}
&=EU\ket{E}\\
&=UE\ket{E}\\
&=UH\ket{E}.
\end{align}
Следовательно
\begin{equation}
(HU-UH)\ket{E}=0.
\end{equation}

\begin{classroomqa}
\audiencequestion{Можно ли повторить вывод, начиная с первого соотношения? \lecturetimestamp{00:35:55}}
\lecturerresponse{Начните с \(H\ket{E}=E\ket{E}\). Симметрия означает, что \(U\ket{E}\) имеет ту же энергию, так что \(HU\ket{E}=EU\ket{E}\). Перенесите скаляр \(E\) через \(U\), замените \(E\ket{E}\) на \(H\ket{E}\) и вычтите. Получается \((HU-UH)\ket{E}=0\).}
\end{classroomqa}

\begin{figure}[H]
\centering
\includegraphics[width=0.70\textwidth]{lecture_04_figure_05.png}
\lectureframe{00:38:51.066}
\caption{Полный вывод на доске. Изогнутые вспомогательные пометки связывают утверждение об одной и той же энергии с подстановкой \(E\ket{E}=H\ket{E}\); последняя строка — \((HU-UH)\ket{E}=0\).}
\end{figure}

Этот аргумент верен для каждого собственного состояния энергии. Предполагая, что эти состояния образуют полный базис, оператор, обращающий каждое из них в нуль, есть нулевой оператор:
\begin{equation}
[H,U]=0.
\end{equation}
Это критерий того, что не зависящее от времени унитарное преобразование является симметрией гамильтониана.

Для бесконечно малого преобразования,
\begin{equation}
U(\epsilon)=\mathbf 1+i\epsilon L_i+O(\epsilon^2),
\end{equation}
оператор тождества и так коммутирует с \(H\), поэтому каждый генератор должен удовлетворять
\begin{equation}
[L_i,H]=0.
\end{equation}

\begin{figure}[H]
\centering
\includegraphics[width=0.47\textwidth]{lecture_04_figure_06.png}
\lectureframe{00:42:14.296}
\caption{Критерий симметрии, сжато записанный на доске: \([U,H]=0\), а затем его форма через генератор, приведённая к стандартному виду по частично размытой нижней подписи: \([L_i,H]=0\).}
\end{figure}

Если \(L_i\) не зависит явно от времени, уравнение Гейзенберга даёт
\begin{equation}
\frac{dL_i}{dt}
=i[H,L_i]=0
\qquad(\hbar=1).
\end{equation}
Таким образом, генераторы симметрии являются сохраняющимися зарядами. Симметрия пространственных переносов даёт закон сохранения импульса; вращательная симметрия даёт закон сохранения момента импульса. \editorialnote{Оговорка об отсутствии явной зависимости от времени стандартна и неявно присутствует в лекционной формулировке закона сохранения.}

\section{Преобразуйте всю физическую систему}

Выражение «поворот является симметрией» можно употребить неправильно, если преобразуется только часть установки. Свободный спин можно повернуть, не изменяя его энергии. Однако если поместить его в фиксированное магнитное поле и повернуть только спин, его энергия Зеемана изменится. Это не опровергает вращательную симметрию; это означает, что поле удерживалось фиксированным, тогда как спин поворачивали относительно него. Если повернуть спин и поле вместе или использовать внешнее поле, которое само инвариантно, энергия не изменится.

Та же проблема возникает в атоме водорода. Смещение только электрона относительно протона изменяет расстояние между ними и, следовательно, кулоновскую энергию. Симметрия переносов действует на полную электрон-протонную систему, смещая оба на одну и ту же величину.

\begin{classroomqa}
\audiencequestion{Должно ли внешнее поле само обладать этой симметрией? \lecturetimestamp{00:43:09}}
\lecturerresponse{Либо фиксированное внешнее поле должно уже быть инвариантным относительно данной операции, либо его нужно преобразовывать вместе с остальной системой. Относительный поворот между спином и неповернутым полем — это другая физическая конфигурация, а не симметрийное преобразование всей установки.}
\end{classroomqa}

Вращательное вырождение в атоме даёт ясный положительный пример. При фиксированном угловом моменте \(l\) повороты смешивают состояния с магнитными квантовыми числами
\begin{equation}
m=-l,-l+1,\ldots,l.
\end{equation}
В отсутствие внешнего магнитного или электрического поля все члены мультиплета имеют одну и ту же энергию:
\begin{equation}
E_{l,m}=E_l.
\end{equation}
Симметрия связывает различные состояния и обусловливает их вырождение.

\section{Что меняет суперсимметрия}

Все знакомые симметрии, обсуждавшиеся до сих пор, сохраняют принадлежность к бозонам или фермионам. Поворот переводит электрон в другой электрон, а не в фотон. Цветовые повороты смешивают три цвета кварка, но не превращают кварк в глюон. Обычные внутренние и пространственно-временные симметрии могут смешивать компоненты или квантовые метки, однако их генераторы бозонны и не меняют фермионную чётность.

Суперсимметрия иная:
\begin{equation}
Q\ket{B}\sim\ket{F},
\qquad
Q\ket{F}\sim\ket{B}.
\end{equation}
Её мультиплеты пересекают границу между бозонами и фермионами. В точной суперсимметричной теории партнёры в одном мультиплете имеют одинаковую массу. Именно такое спаривание требовалось в предыдущих лекциях: согласованные бозонные и фермионные петли могут взаимно компенсировать большие ультрафиолетовые поправки к массам скаляров. Эту конструкцию трудно наглядно представить, поскольку никакое обычное пространственное движение не превращает состояние с целым спином в состояние с полуцелым спином. Поэтому далее мы будем действовать алгебраически.

\section{Прототип суперзаряда}

Пусть \(A,A^\dagger\) уничтожают и создают бозон, а \(C,C^\dagger\) уничтожают и создают фермион. Простейший эрмитов оператор, который взаимно переводит их одночастичные состояния друг в друга, имеет вид
\begin{equation}
Q=C^\dagger A+A^\dagger C.
\end{equation}
Первый член уничтожает бозон и создаёт фермион; второй уничтожает фермион и создаёт бозон. Если \(Q\) действует на однобозонное состояние, член \(A^\dagger C\) обращается в нуль, потому что нет фермиона, который можно было бы уничтожить, тогда как \(C^\dagger A\) осуществляет превращение. На однофермионном состоянии роли меняются местами.

\begin{classroomqa}
\audiencequestion{Создаёт ли эта конструкция одну частицу, являющуюся суперпозицией бозона и фермиона? \lecturetimestamp{00:52:37}}
\lecturerresponse{\(Q\) — оператор, а не состояние. Действуя на чисто бозонное состояние, он порождает фермионное состояние, а действуя на чисто фермионное состояние, он порождает бозонное состояние. Линейно действуя на уже произвольную суперпозицию, он порождает другую суперпозицию, как и любой линейный оператор.}
\end{classroomqa}

Затем в аудитории проверяют вакуум. Поскольку оба оператора уничтожения обращают его в нуль,
\begin{equation}
A\ket{0}=0,\qquad C\ket{0}=0,
\end{equation}
имеем
\begin{equation}
Q\ket{0}=0.
\end{equation}
Результат — нулевой вектор, а не вторая копия вакуума.

Этот прототип намеренно схематичен. Бозон со спином 0 может отображаться в любую из компонент фермиона со спином 1/2, поэтому самому генератору нужен спинорный индекс:
\begin{equation}
Q_\alpha\ket{B}\sim\ket{F_\alpha}.
\end{equation}
В упрощённой двухкомпонентной схеме лекции
\begin{equation}
Q_\alpha\ket{F_\beta}\sim\delta_{\alpha\beta}\ket{B}.
\end{equation}

\begin{classroomqa}
\audiencequestion{Как размеры пространств состояний могут совпадать, если бозон со спином 0 имеет одну спиновую компоненту, а фермион со спином 1/2 — две? \lecturetimestamp{00:55:02}}
\lecturerresponse{Одного неиндексированного \(Q\) недостаточно. Есть спинорные компоненты \(Q_\alpha\), каждая из которых способна породить соответствующую спиновую компоненту фермиона. Сам суперзаряд несёт половину единицы спина, потому что меняет спин на половину единицы.}
\end{classroomqa}

Операторы \(Q_\alpha\) являются фермионными, или, точнее, грассманово-нечётными. Каждый входящий в них член содержит нечётное число фермионных операторов, поэтому они меняют фермионную чётность. Обычные измеримые наблюдаемые чётны. Следует различать нечётный суперзаряд и нечётный параметр преобразования: конечное суперсимметричное преобразование строится из чётной комбинации, такой как \(\epsilon^\alpha Q_\alpha+\bar\epsilon_{\dot\alpha}\bar Q^{\dot\alpha}\), где и \(\epsilon\), и \(Q\) нечётны. \editorialnote{В лекции генераторы неформально называются числами Грассмана. Точнее, суперзаряды — это нечётные операторы, тогда как антикоммутирующие параметры — числа Грассмана; их произведение чётно и может входить в конечное преобразование.}

\section{Бозонные и фермионные алгебры операторов}

Для бозонных мод, помеченных \(i,j\),
\begin{align}
[A_i,A_j]&=0,&
[A_i^\dagger,A_j^\dagger]&=0,&
[A_i,A_j^\dagger]&=\delta_{ij}.
\end{align}
Метки могут обозначать импульс, положение, спин или любой полный одночастичный базис. Поскольку операторы рождения коммутируют,
\begin{equation}
A_i^\dagger A_j^\dagger\ket{0}
=A_j^\dagger A_i^\dagger\ket{0}.
\end{equation}
Двухбозонное состояние симметрично относительно перестановки.

Фермионные операторы удовлетворяют антикоммутационным соотношениям:
\begin{align}
\{C_i,C_j\}&=0,&
\{C_i^\dagger,C_j^\dagger\}&=0,&
\{C_i,C_j^\dagger\}&=\delta_{ij},
\end{align}
где
\begin{equation}
\{X,Y\}=XY+YX.
\end{equation}
Для одной моды
\begin{equation}
\{C^\dagger,C^\dagger\}=2(C^\dagger)^2=0,
\end{equation}
поэтому
\begin{equation}
(C^\dagger)^2=0.
\end{equation}
Сам оператор не равен нулю; его квадрат равен нулю. Он может создать один фермион в этой моде, но второе применение даёт нуль. Это принцип запрета Паули на языке операторов.

Для различных меток
\begin{equation}
C_i^\dagger C_j^\dagger\ket{0}
=-C_j^\dagger C_i^\dagger\ket{0},
\end{equation}
что выражает антисимметрию двухфермионного состояния. Алгебраическая замена
\begin{equation}
\begin{gathered}
\text{bosons: commutators}\\
\Updownarrow\\
\text{fermions: anticommutators}
\end{gathered}
\end{equation}
кодирует статистическое различие.

В одной моде может находиться много бозонов. Поэтому большое число фотонов в моде может вести себя как классическое электромагнитное поле с обычной числовой амплитудой. Фермионы не могут накапливаться в одной моде, и нечётное фермионное поле не приобретает обычного коммутирующего классического значения. Это мотивирует введение иной арифметики для учёта фермионов.

\section{Грассмановы числа и чётность}

Обычные числа коммутируют:
\begin{equation}
\alpha_i\alpha_j=\alpha_j\alpha_i.
\end{equation}
Грассмановы генераторы антикоммутируют:
\begin{equation}
\theta_i\theta_j=-\theta_j\theta_i.
\end{equation}
Если положить \(i=j\), то в числовой системе, где \(2\neq0\), получаем
\begin{equation}
\theta_i^2=0.
\end{equation}
Никакое обычное числовое значение, такое как шесть или семь, нельзя приписать \(\theta_i\). Эти символы образуют алгебру, предназначенную для учёта фермионных знаков.

Обычные коэффициенты коммутируют с грассмановыми генераторами:
\begin{equation}
\alpha\theta_i=\theta_i\alpha.
\end{equation}

\begin{classroomqa}
\audiencequestion{Какой объект получается при умножении обычного числа на грассманово число? \lecturetimestamp{01:12:32}}
\lecturerresponse{Оно остаётся нечётным грассмановым элементом. Умножение на обычный коэффициент меняет его масштаб, а не чётность. Линейная оболочка генераторов образует векторное пространство над обычными числами.}
\end{classroomqa}

Алгебра обладает \(\mathbb Z_2\)-чётностью. Обычные коэффициенты и произведения чётного числа генераторов являются чётными; генератор, как и произведение нечётного числа генераторов, нечётен:
\begin{equation}
\begin{aligned}
\text{even}\cdot\text{even}&=\text{even},\\
\text{even}\cdot\text{odd}&=\text{odd},\\
\text{odd}\cdot\text{odd}&=\text{even}.
\end{aligned}
\end{equation}

\begin{classroomqa}
\audiencequestion{Аналогично ли это правилам умножения действительных и мнимых чисел или это обобщение, подобное введению \(i\)? \lecturetimestamp{01:14:25}}
\lecturerresponse{Это обобщение арифметики, но не следует слишком далеко заходить в этой аналогии. Мнимые числа по-прежнему коммутируют и удовлетворяют \(i^2=-1\); нечётные грассмановы величины антикоммутируют и в квадрате дают нуль. Общая схема умножения по чётности не делает эти алгебры одинаковыми.}
\end{classroomqa}

Произведение, такое как \(\theta_1\theta_2\), чётно и коммутирует с каждым однородным элементом алгебры Грассмана, но не является обычным числом. Оно ненулевое, тогда как
\begin{equation}
(\theta_1\theta_2)^2=0.
\end{equation}

\begin{classroomqa}
\audiencequestion{Если \(\theta_1\theta_2\) чётно, почему это не просто обычное число? \lecturetimestamp{01:17:36}}
\lecturerresponse{Чётность означает, что элемент коммутирует в градуированной алгебре; это не означает, что он является действительным или комплексным скаляром. Его нильпотентность сразу отличает его от ненулевого обычного числа.}
\end{classroomqa}

\section{Функции грассмановых переменных}

При одном генераторе любой многочлен обрывается на члене первого порядка:
\begin{equation}
f(\theta)=\alpha+\beta\theta.
\end{equation}
Все степени \(\theta^n\) при \(n\geq2\) обращаются в нуль. Для двух генераторов наиболее общий элемент имеет вид
\begin{equation}
f(\theta_1,\theta_2)
=\alpha+\beta_1\theta_1+\beta_2\theta_2
+\gamma\theta_1\theta_2.
\end{equation}
Если записать последний моном в обратном порядке, его знак меняется:
\begin{equation}
\theta_2\theta_1=-\theta_1\theta_2.
\end{equation}

\begin{classroomqa}
\audiencequestion{Порождает ли умножение \(\theta_1\) на коэффициент \(\beta\) третью независимую грассманову переменную? \lecturetimestamp{01:18:39}}
\lecturerresponse{Нет. Оно создаёт ещё один вектор в линейной оболочке \(\theta_1\), а не новый базисный генератор. При двух независимых генераторах любой нечётный линейный член имеет вид \(\beta_1\theta_1+\beta_2\theta_2\).}
\end{classroomqa}

Для трёх генераторов упорядоченный базис таков:
\begin{equation}
1,\quad
\theta_1,\theta_2,\theta_3,\quad
\theta_1\theta_2,\theta_1\theta_3,\theta_2\theta_3,\quad
\theta_1\theta_2\theta_3.
\end{equation}
В общем случае \(n\) генераторов дают \(2^n\) упорядоченных базисных мономов. Любая аналитическая функция от конечного числа грассмановых генераторов обрывается на таком многочлене, поскольку повторяющиеся генераторы обращаются в нуль. \editorialnote{Этот \(2^n\)-мерный базис явно задаёт схему, которую лекция строит для одного, двух и трёх генераторов.}

\section{Дифференцирование по грассмановым переменным}

Для одной переменной определим
\begin{equation}
\frac{\partial}{\partial\theta}\alpha=0,
\qquad
\frac{\partial}{\partial\theta}\theta=1,
\end{equation}
так что
\begin{equation}
\frac{\partial}{\partial\theta}
(\alpha+\beta\theta)=\beta.
\end{equation}
Для двух переменных выберем левые производные и будем сохранять канонический порядок \(\theta_1\theta_2\):
\begin{align}
\frac{\partial^L f}{\partial\theta_1}
&=\beta_1+\gamma\theta_2,\\
\frac{\partial^L f}{\partial\theta_2}
&=\beta_2-\gamma\theta_1.
\end{align}
Знак минус во второй строке возникает потому, что нечётная производная должна пройти через \(\theta_1\), прежде чем продифференцировать \(\theta_2\).

Это выражается градуированным правилом произведения. Если \(A\) имеет чётность \(|A|\in\{0,1\}\), то
\begin{equation}
\frac{\partial^L}{\partial\theta}(AB)
=\left(\frac{\partial^L A}{\partial\theta}\right)B
+(-1)^{|A|}A
\left(\frac{\partial^L B}{\partial\theta}\right).
\end{equation}
По духу это обычное правило, но в нём присутствует знак, требуемый антикоммутацией. Интегрирование по грассмановым переменным также можно определить и, поскольку пространство функций конечно, оно сводится к короткой таблице. Лекция откладывает эту таблицу.

\section{Разрешение загадки с экспонентой}

С одним грассмановым генератором и обычным коэффициентом \(a\)
\begin{equation}
e^{a\theta}
=1+a\theta,
\end{equation}
поскольку все высшие степени обращаются в нуль. Для двух экспонент с одним и тем же генератором выполняется
\begin{align}
e^{a\theta}e^{b\theta}
&=(1+a\theta)(1+b\theta)\\
&=1+(a+b)\theta\\
&=e^{(a+b)\theta}.
\end{align}

Затем аудитория рассматривает различные генераторы. Непосредственное умножение даёт
\begin{equation}
e^{\theta_1}e^{\theta_2}
=(1+\theta_1)(1+\theta_2)
=1+\theta_1+\theta_2+\theta_1\theta_2.
\end{equation}
Поначалу разложение на доске наводит на мысль, что квадрат \(\theta_1+\theta_2\) может дать последний член. Один из слушателей замечает знак:
\begin{equation}
(\theta_1+\theta_2)^2
=\theta_1^2+\theta_1\theta_2
 +\theta_2\theta_1+\theta_2^2
=0.
\end{equation}

\begin{classroomqa}
\audiencequestion{Квадратичный член сокращается потому, что он содержит \(\theta_1\theta_2+\theta_2\theta_1\)? \lecturetimestamp{01:24:00}}
\lecturerresponse{Да. Эта поправка меняет вывод: \(e^{\theta_1+\theta_2}=1+\theta_1+\theta_2\), тогда как упорядоченное произведение \(e^{\theta_1}e^{\theta_2}\) тоже содержит \(\theta_1\theta_2\). Наивный закон сложения экспонент не применим к различным антикоммутирующим аргументам.}
\end{classroomqa}

Противоречия нет. Даже для обычных матриц \(e^Ae^B=e^{A+B}\) выполняется только тогда, когда \(A\) и \(B\) коммутируют. Здесь
\begin{equation}
[\theta_1,\theta_2]
=\theta_1\theta_2-\theta_2\theta_1
=2\theta_1\theta_2.
\end{equation}
Этот коммутатор чётен и коммутирует с обоими генераторами, поэтому ряд Бейкера—Кэмпбелла—Хаусдорфа обрывается:
\begin{equation}
e^{\theta_1}e^{\theta_2}
=\exp\left(
\theta_1+\theta_2+\frac12[\theta_1,\theta_2]
\right)
=e^{\theta_1+\theta_2+\theta_1\theta_2}.
\end{equation}
Поскольку аргумент экспоненты справа нильпотентен, его разложение в точности равно
\(1+\theta_1+\theta_2+\theta_1\theta_2\). \editorialnote{Этот расчёт по формуле Бейкера—Кэмпбелла—Хаусдорфа даёт то разрешение, которое лекция в аудитории оставляет в виде вопроса. Лекция правильно обнаруживает несостоятельность наивного закона, но заканчивается прежде, чем выявляется недостающий коммутаторный член.}

\begin{classroomqa}
\audiencequestion{Что это за число \(1+\theta_1\), и можно ли складывать обычную и грассманову части? \lecturetimestamp{01:26:16}}
\lecturerresponse{Это неоднородный элемент алгебры Грассмана с чётной частью \(1\) и нечётной частью \(\theta_1\). Такие суммы допустимы, но они не являются ни обычными скалярами, ни однородными нечётными элементами.}
\end{classroomqa}

Ближе к концу проверяется ещё один кажущийся парадокс:
\begin{equation}
(\theta_1+\theta_2)^2=0.
\end{equation}
Это не влечёт \(\theta_1\theta_2=0\), потому что два перекрёстных члена сокращаются, а не обращаются в нуль по отдельности.

\begin{classroomqa}
\audiencequestion{Если квадрат всякого нечётного грассманова элемента равен нулю, означает ли это, что произведение любых двух грассмановых генераторов равно нулю? \lecturetimestamp{01:32:19}}
\lecturerresponse{Нет. Равенство \((\theta_1+\theta_2)^2=0\) означает \(\theta_1\theta_2+\theta_2\theta_1=0\). Оно доказывает антикоммутацию, а не то, что какое-либо из двух упорядоченных произведений обращается в нуль. Произведение различных генераторов может быть ненулевым.}
\end{classroomqa}

\begin{classroomqa}
\audiencequestion{Предполагаем ли мы обычное определение через степенной ряд, когда записываем эти экспоненты? \lecturetimestamp{01:34:42}}
\lecturerresponse{Да, и это определение согласовано, потому что каждый ряд обрывается. Ломается не сама экспонента, а неоправданное тождество \(e^Ae^B=e^{A+B}\) для некоммутирующих \(A\) и \(B\). Поправка через коммутатор восстанавливает правильный закон композиции.}
\end{classroomqa}

Такой исследовательский финал уместен. Обычные симметрии привели нас от унитарных преобразований к генераторам и коммутаторам; суперсимметрия привела нас от нечётных генераторов к антикоммутирующим переменным. В обоих случаях центральным объектом является закон умножения. Для вращений он проявляется в \([L_i,L_j]\); для грассмановых экспонент — в коммутаторном члене, который препятствует наивному сложению. Эта алгебра непривычна, но она непротиворечива, конечна и в точности подходит для фермионных преобразований, которых требует суперсимметрия.